\section{Introduction}
\label{sec:introduction}

Choosing between two similarly valued snacks should take longer than choosing
between a favorite and a much less valued alternative. The drift-diffusion
model (DDM) links this prediction to the probability of choosing each good.
When values are close, evidence drifts slowly toward either decision threshold;
with a larger value difference, choices become more predictable and quicker on
average. I use the formulation presented by
\textcite[slides 10--11 and 24--26]{woodford2026ddm} in Lecture 5.

The individual file, \texttt{indv-Data.csv}, records 435 choices under my UNI,
\texttt{ts3686}. The pooled file,
\texttt{popData.csv}, supplies 15 class bins.
The first good in each individual record is treated as the left good and the
second as the right good. Their mean bids from Section 1 of the experiment
give $v_L$ and $v_R$, so the value difference is $\Delta=v_R-v_L$.
There are 368 distinct pairs of goods, some of which I encountered more than
once; each choice is a separate trial.
